\hypertarget{variables-expressions-et-uxe9noncuxe9s}{%
\section{Variables, expressions et
énoncés}\label{variables-expressions-et-uxe9noncuxe9s}}

\hypertarget{valeurs-et-types}{%
\subsection{Valeurs et types}\label{valeurs-et-types}}

Une valeur est l\textquotesingle une des choses de base avec lesquelles
un programme fonctionne, comme une lettre ou un numéro. Les valeurs que
nous avons vues jusqu\textquotesingle ici sont 1, 2, et
\textquotesingle Hello, World!\textquotesingle.

Ces valeurs appartiennent à différents types: \texttt{2} est un entier,
et \texttt{"Hello,\ World!"} est une chaîne de caractère
(\emph{string}). Vous (et l\textquotesingle interprète) pouvez
identifier ces string parce qu\textquotesingle ils sont placées entre
des guillemets (\texttt{"}).

Si vous ne savez pas quel type de valeur a une valeur,
l\textquotesingle interprète peut vous aider.

\begin{Shaded}
\begin{Highlighting}[]
\OperatorTok{\textgreater{}\textgreater{}\textgreater{}} \BuiltInTok{type}\NormalTok{(}\StringTok{"Hello World!"}\NormalTok{)}
\OperatorTok{\textless{}}\KeywordTok{class} \StringTok{\textquotesingle{}str\textquotesingle{}}\OperatorTok{\textgreater{}}
\OperatorTok{\textgreater{}\textgreater{}\textgreater{}} \BuiltInTok{type}\NormalTok{(}\DecValTok{17}\NormalTok{)}
\OperatorTok{\textless{}}\KeywordTok{class} \StringTok{\textquotesingle{}int\textquotesingle{}}\OperatorTok{\textgreater{}}
\OperatorTok{\textgreater{}\textgreater{}\textgreater{}} \BuiltInTok{type}\NormalTok{(}\FloatTok{1.3}\NormalTok{)}
\OperatorTok{\textless{}}\KeywordTok{class} \StringTok{\textquotesingle{}float\textquotesingle{}}\OperatorTok{\textgreater{}}
\OperatorTok{\textgreater{}\textgreater{}\textgreater{}} \BuiltInTok{type}\NormalTok{(}\StringTok{\textquotesingle{}17\textquotesingle{}}\NormalTok{)}
\OperatorTok{\textless{}}\KeywordTok{class} \StringTok{\textquotesingle{}str\textquotesingle{}}\OperatorTok{\textgreater{}}
\OperatorTok{\textgreater{}\textgreater{}\textgreater{}} \BuiltInTok{type}\NormalTok{(}\StringTok{\textquotesingle{}3.2\textquotesingle{}}\NormalTok{)}
\OperatorTok{\textless{}}\KeywordTok{class} \StringTok{\textquotesingle{}str\textquotesingle{}}\OperatorTok{\textgreater{}}
\end{Highlighting}
\end{Shaded}

\hypertarget{variables}{%
\subsection{Variables}\label{variables}}

L\textquotesingle une des fonctionnalités les plus puissantes
d\textquotesingle un langage de programmation est la capacité de
manipuler des variables. Une variable est un nom qui fait référence à
une valeur.

Une instruction d\textquotesingle affectation (\emph{assignment
statement}) crée de nouvelles variables et leur donne des valeurs :

\begin{Shaded}
\begin{Highlighting}[]
\OperatorTok{\textgreater{}\textgreater{}\textgreater{}}\NormalTok{ message }\OperatorTok{=} \StringTok{"Voici quelque chose de différent"}
\OperatorTok{\textgreater{}\textgreater{}\textgreater{}}\NormalTok{ n }\OperatorTok{=} \DecValTok{17}
\OperatorTok{\textgreater{}\textgreater{}\textgreater{}}\NormalTok{ pi }\OperatorTok{=} \FloatTok{3.1415926535897932}
\end{Highlighting}
\end{Shaded}

Cet exemple effectue trois affectations. La première assigne une chaîne
à une nouvelle variable nommée \texttt{message} ; la deuxième donne
l\textquotesingle entier \texttt{17} à \texttt{n} ; la troisième assigne
la valeur (approximative) de π à \texttt{pi}.

Le type d\textquotesingle une variable est le type de la valeur à
laquelle elle fait référence.

\begin{Shaded}
\begin{Highlighting}[]
\OperatorTok{\textgreater{}\textgreater{}\textgreater{}} \BuiltInTok{type}\NormalTok{(message)}
\OperatorTok{\textless{}}\KeywordTok{class} \StringTok{\textquotesingle{}str\textquotesingle{}}\OperatorTok{\textgreater{}}
\OperatorTok{\textgreater{}\textgreater{}\textgreater{}} \BuiltInTok{type}\NormalTok{(n)}
\OperatorTok{\textless{}}\KeywordTok{class} \StringTok{\textquotesingle{}int\textquotesingle{}}\OperatorTok{\textgreater{}}
\OperatorTok{\textgreater{}\textgreater{}\textgreater{}} \BuiltInTok{type}\NormalTok{(pi)}
\OperatorTok{\textless{}}\KeywordTok{class} \StringTok{\textquotesingle{}float\textquotesingle{}}\OperatorTok{\textgreater{}}
\end{Highlighting}
\end{Shaded}

\hypertarget{noms-de-variables-et-mots-cluxe9s}{%
\subsection{Noms de variables et
mots-clés}\label{noms-de-variables-et-mots-cluxe9s}}

Les programmeurs choisissent généralement des noms significatifs pour
leurs variables. Ils documentent à quoi sert la variable.

Les noms de variables peuvent être aussi longs que vous le souhaitez.
Ils peuvent contenir des lettres et des chiffres, mais ils doivent
commencer par une lettre. Il est légal d\textquotesingle utiliser des
majuscules, mais c\textquotesingle est une bonne idée de commencer les
noms de variables par une minuscule (vous verrez pourquoi plus tard).

Le caractère de soulignement (\emph{underscore}), \texttt{\_}, peut
apparaître dans un nom. Il est souvent utilisé dans les noms composés de
plusieurs mots, comme \texttt{mon\_nom} ou
\texttt{vitesse\_hirondelle\_a\_vide}.

Si vous donnez à une variable un nom illégal, vous obtenez une erreur de
syntaxe :

\begin{Shaded}
\begin{Highlighting}[]
\NormalTok{\textgreater{}\textgreater{}\textgreater{} 76trombones = \textquotesingle{}big parade\textquotesingle{}}
\NormalTok{File "\textless{}stdin\textgreater{}", line 1}
\NormalTok{    76trombones = \textquotesingle{}big parade\textquotesingle{}}
\NormalTok{    \^{}}
\NormalTok{SyntaxError: invalid decimal literal}
\NormalTok{\textgreater{}\textgreater{}\textgreater{} more@ = 1000000}
\NormalTok{  File "\textless{}stdin\textgreater{}", line 1}
\NormalTok{    more@ = 1000000}
\NormalTok{          \^{}}
\NormalTok{SyntaxError: invalid syntax}
\NormalTok{\textgreater{}\textgreater{}\textgreater{} class = \textquotesingle{}Zymologie théorique avancée\textquotesingle{}}
\NormalTok{    File "\textless{}stdin\textgreater{}", line 1}
\NormalTok{    class = \textquotesingle{}Zymologie théorique avancée\textquotesingle{}}
\NormalTok{            \^{}}
\NormalTok{SyntaxError: invalid syntax}
\end{Highlighting}
\end{Shaded}

\texttt{76trombones} est illégal car il ne commence pas par une lettre.
\texttt{more@} est illégal car il contient un caractère illégal,
\texttt{@}. Mais quel est le problème avec \texttt{class} ?

Il s\textquotesingle avère que \texttt{class} est l\textquotesingle un
des mots-clés de Python. L\textquotesingle interpréteur utilise les
mots-clés pour reconnaître la structure du programme, et ils ne peuvent
pas être utilisés comme noms de variables.

Python 3.14 possède 31 mots-clés :

\texttt{False}, \texttt{None}, \texttt{True}, \texttt{and}, \texttt{as},
\texttt{assert}, \texttt{async}, \texttt{await}, \texttt{break},
\texttt{class}, \texttt{continue}, \texttt{def}, \texttt{del},
\texttt{elif}, \texttt{else}, \texttt{except}, \texttt{finally},
\texttt{for}, \texttt{from}, \texttt{global}, \texttt{if},
\texttt{import}, \texttt{in}, \texttt{is}, \texttt{lambda},
\texttt{nonlocal}, \texttt{not}, \texttt{or}, \texttt{pass},
\texttt{raise}, \texttt{return}, \texttt{try}, \texttt{while},
\texttt{with} et \texttt{yield}.

Vous voudrez peut-être garder cette liste à portée de main. Si
l\textquotesingle interpréteur se plaint d\textquotesingle un de vos
noms de variables et que vous ne savez pas pourquoi, vérifiez
s\textquotesingle il figure dans cette liste. Vous pouvez aussi
simplementdemander à l\textquotesingle interpréteur de vous la donner:

\begin{Shaded}
\begin{Highlighting}[]
\OperatorTok{\textgreater{}\textgreater{}\textgreater{}} \ImportTok{import}\NormalTok{ keyword}
\OperatorTok{\textgreater{}\textgreater{}\textgreater{}} \BuiltInTok{print}\NormalTok{(keyword.kwlist)}
\NormalTok{[}\StringTok{\textquotesingle{}False\textquotesingle{}}\NormalTok{, }\StringTok{\textquotesingle{}None\textquotesingle{}}\NormalTok{, }\StringTok{\textquotesingle{}True\textquotesingle{}}\NormalTok{, }\StringTok{\textquotesingle{}and\textquotesingle{}}\NormalTok{, }\StringTok{\textquotesingle{}as\textquotesingle{}}\NormalTok{, }\StringTok{\textquotesingle{}assert\textquotesingle{}}\NormalTok{, }\StringTok{\textquotesingle{}async\textquotesingle{}}\NormalTok{, }\StringTok{\textquotesingle{}await\textquotesingle{}}\NormalTok{, }\StringTok{\textquotesingle{}break\textquotesingle{}}\NormalTok{, }\StringTok{\textquotesingle{}class\textquotesingle{}}\NormalTok{, }\StringTok{\textquotesingle{}continue\textquotesingle{}}\NormalTok{, }\StringTok{\textquotesingle{}def\textquotesingle{}}\NormalTok{, }\StringTok{\textquotesingle{}del\textquotesingle{}}\NormalTok{, }\StringTok{\textquotesingle{}elif\textquotesingle{}}\NormalTok{, }\StringTok{\textquotesingle{}else\textquotesingle{}}\NormalTok{, }\StringTok{\textquotesingle{}except\textquotesingle{}}\NormalTok{, }\StringTok{\textquotesingle{}finally\textquotesingle{}}\NormalTok{, }\StringTok{\textquotesingle{}for\textquotesingle{}}\NormalTok{, }\StringTok{\textquotesingle{}from\textquotesingle{}}\NormalTok{, }\StringTok{\textquotesingle{}global\textquotesingle{}}\NormalTok{, }\StringTok{\textquotesingle{}if\textquotesingle{}}\NormalTok{, }\StringTok{\textquotesingle{}import\textquotesingle{}}\NormalTok{, }\StringTok{\textquotesingle{}in\textquotesingle{}}\NormalTok{, }\StringTok{\textquotesingle{}is\textquotesingle{}}\NormalTok{, }\StringTok{\textquotesingle{}lambda\textquotesingle{}}\NormalTok{, }\StringTok{\textquotesingle{}nonlocal\textquotesingle{}}\NormalTok{, }\StringTok{\textquotesingle{}not\textquotesingle{}}\NormalTok{, }\StringTok{\textquotesingle{}or\textquotesingle{}}\NormalTok{, }\StringTok{\textquotesingle{}pass\textquotesingle{}}\NormalTok{, }\StringTok{\textquotesingle{}raise\textquotesingle{}}\NormalTok{, }\StringTok{\textquotesingle{}return\textquotesingle{}}\NormalTok{, }\StringTok{\textquotesingle{}try\textquotesingle{}}\NormalTok{, }\StringTok{\textquotesingle{}while\textquotesingle{}}\NormalTok{, }\StringTok{\textquotesingle{}with\textquotesingle{}}\NormalTok{, }\StringTok{\textquotesingle{}yield\textquotesingle{}}\NormalTok{]}
\end{Highlighting}
\end{Shaded}

\hypertarget{opuxe9rateurs-et-opuxe9randes}{%
\subsection{Opérateurs et
opérandes}\label{opuxe9rateurs-et-opuxe9randes}}

Les opérateurs sont des symboles spéciaux qui représentent des calculs
comme l\textquotesingle addition et la multiplication. Les valeurs
auxquelles l\textquotesingle opérateur est appliqué
s\textquotesingle appellent des opérandes.

Les opérateurs \texttt{+}, \texttt{-}, \texttt{*}, \texttt{/},
\texttt{//} et \texttt{**} effectuent l\textquotesingle addition, la
soustraction, la multiplication, la division, la division entière, et
l\textquotesingle exponentiation, comme dans les exemples suivants :

\begin{Shaded}
\begin{Highlighting}[]
\DecValTok{20}\OperatorTok{+}\DecValTok{32}
\NormalTok{hour}\OperatorTok{{-}}\DecValTok{1}
\NormalTok{hour}\OperatorTok{*}\DecValTok{60}\OperatorTok{+}\NormalTok{minute}
\NormalTok{minute}\OperatorTok{/}\DecValTok{60}
\DecValTok{5}\OperatorTok{**}\DecValTok{2}
\NormalTok{(}\DecValTok{5}\OperatorTok{+}\DecValTok{9}\NormalTok{)}\OperatorTok{*}\NormalTok{(}\DecValTok{15}\OperatorTok{{-}}\DecValTok{7}\NormalTok{)}
\end{Highlighting}
\end{Shaded}

Dans d\textquotesingle autres langages, \texttt{\^{}} est utilisé pour
l\textquotesingle exponentiation, mais en Python, c\textquotesingle est
un opérateur bit à bit (XOR).

\hypertarget{expressions-et-instructions}{%
\subsection{Expressions et
instructions}\label{expressions-et-instructions}}

Une expression est une combinaison de valeurs, de variables et
d\textquotesingle opérateurs. Une valeur seule est considérée comme une
expression, de même qu\textquotesingle une variable, donc ce qui suit
sont toutes des expressions légales :

\begin{Shaded}
\begin{Highlighting}[]
\DecValTok{17}
\NormalTok{x}
\NormalTok{x }\OperatorTok{+} \DecValTok{17}
\end{Highlighting}
\end{Shaded}

Une instruction (\emph{statement}) est une unité de code que
l\textquotesingle interpréteur Python peut exécuter.

Techniquement, une expression est aussi une instruction, mais il est
probablement plus simple de les considérer comme des choses différentes.
La différence importante est qu\textquotesingle une expression a une
valeur ; une instruction n\textquotesingle en a pas.

\hypertarget{mode-interactif-et-mode-script}{%
\subsection{Mode interactif et mode
script}\label{mode-interactif-et-mode-script}}

L\textquotesingle un des avantages de travailler avec un langage
interprété est que vous pouvez tester des morceaux de code en mode
interactif avant de les mettre dans un script.

Par exemple, si vous utilisez Python comme calculatrice, vous pourriez
taper :

\begin{Shaded}
\begin{Highlighting}[]
\OperatorTok{\textgreater{}\textgreater{}\textgreater{}}\NormalTok{ miles }\OperatorTok{=} \FloatTok{26.2}
\OperatorTok{\textgreater{}\textgreater{}\textgreater{}}\NormalTok{ miles }\OperatorTok{*} \FloatTok{1.61}
\FloatTok{42.182}
\end{Highlighting}
\end{Shaded}

La première ligne assigne une valeur à \texttt{miles}, mais
n\textquotesingle a pas d\textquotesingle effet visible. La deuxième
ligne est une expression, donc l\textquotesingle interpréteur
l\textquotesingle évalue et affiche le résultat. Nous apprenons donc
qu\textquotesingle un marathon fait environ 42 kilomètres.

Mais si vous tapez le même code dans un script et
l\textquotesingle exécutez, vous n\textquotesingle obtenez aucune sortie
du tout. En mode script, une expression seule n\textquotesingle a aucun
effet visible. Python évalue l\textquotesingle expression, mais il
n\textquotesingle affiche pas la valeur à moins qu\textquotesingle on ne
lui demande :

\begin{Shaded}
\begin{Highlighting}[]
\NormalTok{miles }\OperatorTok{=} \FloatTok{26.2}
\BuiltInTok{print}\NormalTok{(miles }\OperatorTok{*} \FloatTok{1.61}\NormalTok{)}
\end{Highlighting}
\end{Shaded}

Ce comportement peut être déroutant au début.

Un script contient généralement une séquence
d\textquotesingle instructions. S\textquotesingle il y a plus
d\textquotesingle une instruction, les résultats apparaissent un par un
au fur et à mesure de l\textquotesingle exécution.

Par exemple, le script :

\begin{Shaded}
\begin{Highlighting}[]
\BuiltInTok{print}\NormalTok{(}\DecValTok{1}\NormalTok{)}
\NormalTok{x }\OperatorTok{=} \DecValTok{2}
\BuiltInTok{print}\NormalTok{(x)}
\end{Highlighting}
\end{Shaded}

produit la sortie :

\begin{Shaded}
\begin{Highlighting}[]
\NormalTok{1}
\NormalTok{2}
\end{Highlighting}
\end{Shaded}

L\textquotesingle instruction d\textquotesingle affectation ne produit
aucune sortie.

\textbf{Exercice 1}

Tapez les instructions suivantes dans l\textquotesingle interpréteur
Python pour voir ce qu\textquotesingle elles font :

\begin{Shaded}
\begin{Highlighting}[]
\DecValTok{5}
\NormalTok{x }\OperatorTok{=} \DecValTok{5}
\NormalTok{x }\OperatorTok{+} \DecValTok{1}
\end{Highlighting}
\end{Shaded}

Maintenant, mettez ces mêmes instructions dans un script et exécutez-le.
Quelle est la sortie ? Modifiez le script en transformant chaque
expression en un appel à la fonction \texttt{print}, puis exécutez-le à
nouveau.

\hypertarget{ordre-des-opuxe9rations}{%
\subsection{Ordre des opérations}\label{ordre-des-opuxe9rations}}

Lorsqu\textquotesingle il y a plus d\textquotesingle un opérateur dans
une expression, l\textquotesingle ordre d\textquotesingle évaluation
dépend des règles de priorité. Pour les opérateurs mathématiques, Python
suit la convention mathématique (PEMDAS) :

\begin{itemize}
\tightlist
\item
  \textbf{Parenthèses} ont la plus haute priorité et peuvent être
  utilisées pour forcer une expression à s\textquotesingle évaluer dans
  l\textquotesingle ordre de votre choix.
\item
  \textbf{Exponentiation} a la priorité suivante.
\item
  \textbf{Multiplication} et \textbf{Division} ont la même priorité, qui
  est plus élevée que l\textquotesingle{}\textbf{Addition} et la
  \textbf{Soustraction} (qui ont également la même priorité entre
  elles).
\item
  Les opérateurs avec la même priorité sont évalués de gauche à droite
  (à l\textquotesingle exception de l\textquotesingle exponentiation).
\end{itemize}

\hypertarget{opuxe9rations-sur-les-chauxeenes-de-caractuxe8res}{%
\subsection{Opérations sur les chaînes de
caractères}\label{opuxe9rations-sur-les-chauxeenes-de-caractuxe8res}}

En général, on ne peut pas effectuer d\textquotesingle opérations
mathématiques sur les chaînes de caractères, même si elles ressemblent à
des nombres.

L\textquotesingle opérateur \texttt{+} fonctionne avec les chaînes, mais
il effectue une concaténation, ce qui signifie joindre les chaînes en
les liant bout à bout. Par exemple :

\begin{Shaded}
\begin{Highlighting}[]
\NormalTok{first }\OperatorTok{=} \StringTok{\textquotesingle{}Paru\textquotesingle{}}
\NormalTok{second }\OperatorTok{=} \StringTok{\textquotesingle{}line\textquotesingle{}}
\BuiltInTok{print}\NormalTok{(first }\OperatorTok{+}\NormalTok{ second)}
\end{Highlighting}
\end{Shaded}

La sortie de ce programme est \texttt{Paruline}.

L\textquotesingle opérateur \texttt{*} fonctionne également sur les
chaînes ; il effectue la répétition. Par exemple,
\texttt{\textquotesingle{}Spam\textquotesingle{}*3} donne
\texttt{\textquotesingle{}SpamSpamSpam\textquotesingle{}}.

\hypertarget{commentaires}{%
\subsection{Commentaires}\label{commentaires}}

À mesure que les programmes deviennent plus gros et plus complexes, ils
deviennent plus difficiles à lire. Pour cette raison, il est bon
d\textquotesingle ajouter des notes à vos programmes pour expliquer en
langage naturel ce que fait le programme. Ces notes
s\textquotesingle appellent des commentaires, et ils commencent par le
symbole \texttt{\#} :

\begin{Shaded}
\begin{Highlighting}[]
\CommentTok{\# calcule le pourcentage de l\textquotesingle{}heure écoulée}
\NormalTok{percentage }\OperatorTok{=}\NormalTok{ (minute }\OperatorTok{*} \DecValTok{100}\NormalTok{) }\OperatorTok{/} \DecValTok{60}
\end{Highlighting}
\end{Shaded}

Vous pouvez également mettre des commentaires à la fin
d\textquotesingle une ligne :

\begin{Shaded}
\begin{Highlighting}[]
\NormalTok{percentage }\OperatorTok{=}\NormalTok{ (minute }\OperatorTok{*} \DecValTok{100}\NormalTok{) }\OperatorTok{/} \DecValTok{60}  \CommentTok{\# pourcentage d\textquotesingle{}une heure}
\end{Highlighting}
\end{Shaded}

Tout ce qui va du \texttt{\#} jusqu\textquotesingle à la fin de la ligne
est ignoré.

Les commentaires sont très utiles pour expliquer les caractéristiques
non évidentes du code (le \emph{pourquoi} plutôt que le \emph{comment}).

\hypertarget{duxe9bogage}{%
\subsection{Débogage}\label{duxe9bogage}}

À ce stade, l\textquotesingle erreur de syntaxe la plus probable que
vous ferez est un nom de variable illégal, comme
l\textquotesingle utilisation de mots-clés ou de caractères interdits.

L\textquotesingle erreur d\textquotesingle exécution la plus probable
est d\textquotesingle essayer d\textquotesingle utiliser une variable
avant de lui avoir assigné une valeur (une \texttt{NameError}). Cela
peut arriver si vous orthographiez mal un nom de variable.
N\textquotesingle oubliez pas que les variables sont sensibles à la
casse.

L\textquotesingle erreur sémantique la plus fréquente vient souvent de
l\textquotesingle ordre des opérations.

\hypertarget{glossaire}{%
\subsection{Glossaire}\label{glossaire}}

\begin{description}
\item[valeur]
L\textquotesingle une des unités de base des données
qu\textquotesingle un programme manipule.
\item[type]
Une catégorie de valeurs.
\end{description}

entier integer Un type qui représente des nombres entiers.

virgule flottante float Un type qui représente des nombres avec des
parties fractionnaires.

chaîne de caractères string Un type qui représente des séquences de
caractères.

\begin{description}
\item[variable]
Un nom qui fait référence à une valeur.
\end{description}

instruction statement Une section de code qui représente une commande ou
une action.

affectation assignment Une instruction qui assigne une valeur à une
variable.

diagramme d\textquotesingle état state diagram Une représentation
graphique d\textquotesingle un ensemble de variables et des valeurs
auxquelles elles font référence.

mot-clé keyword Un mot réservé utilisé par le compilateur pour analyser
un programme.

opérateur operator Un symbole spécial qui représente un calcul simple.

opérande operand L\textquotesingle une des valeurs sur lesquelles un
opérateur agit.

division entière floor division L\textquotesingle opération qui divise
deux nombres et coupe la partie fractionnaire.

\begin{description}
\item[expression]
Une combinaison de variables, d\textquotesingle opérateurs et de valeurs
qui représente une valeur de résultat unique.
\end{description}

évaluer evaluate Simplifier une expression en effectuant les opérations
pour obtenir une valeur unique.

règles de priorité rules of precedence L\textquotesingle ensemble des
règles régissant l\textquotesingle ordre dans lequel les expressions
sont évaluées.

concaténer concatenate Joindre deux opérandes bout à bout.

commentaire comment Information dans un programme destinée aux lecteurs
du code source.

\hypertarget{exercices}{%
\subsection{Exercices}\label{exercices}}

\textbf{Exercice 2}

Supposons que nous exécutions les instructions
d\textquotesingle affectation suivantes :

\begin{Shaded}
\begin{Highlighting}[]
\NormalTok{width }\OperatorTok{=} \DecValTok{17}
\NormalTok{height }\OperatorTok{=} \FloatTok{12.0}
\NormalTok{delimiter }\OperatorTok{=} \StringTok{\textquotesingle{}.\textquotesingle{}}
\end{Highlighting}
\end{Shaded}

Pour chacune des expressions suivantes, déterminez la valeur de
l\textquotesingle expression et le type (de la valeur de
l\textquotesingle expression) :

\begin{enumerate}
\def\labelenumi{\arabic{enumi}.}
\tightlist
\item
  \texttt{width/2}
\item
  \texttt{width/2.0}
\item
  \texttt{height/3}
\item
  \texttt{1\ +\ 2\ *\ 5}
\item
  \texttt{delimiter\ *\ 5}
\end{enumerate}

\textbf{Exercice 3}

Pratiquez l\textquotesingle utilisation de
l\textquotesingle interpréteur Python comme calculatrice :

\begin{enumerate}
\def\labelenumi{\arabic{enumi}.}
\tightlist
\item
  Le volume d\textquotesingle une sphère de rayon r est 4/3 π r³. Quel
  est le volume d\textquotesingle une sphère de rayon 5 ?
\item
  Supposons que le prix de couverture d\textquotesingle un livre soit de
  24,95 \$, mais que les librairies obtiennent une réduction de 40\%.
  Les frais d\textquotesingle expédition sont de 3 \$ pour le premier
  exemplaire et de 75 cents pour chaque exemplaire supplémentaire. Quel
  est le coût total de gros pour 60 exemplaires ?
\item
  Si je quitte ma maison à 6h52 et que je cours 1 mile à un rythme
  facile (8:15 par mile), puis 3 miles à un rythme soutenu (7:12 par
  mile) et 1 mile à un rythme facile à nouveau, à quelle heure rentré-je
  chez moi pour le petit-déjeuner ?
\end{enumerate}
